\documentclass[oneside,openany,phd]{TMU-Thesis}

% مشخصات پایان‌نامه را در فایلهای faTitle و enTitle وارد نمایید.
% در اینجا برای یکپارچگی، همینجا وارد می‌کنیم

% مشخصات فارسی
\university{تربیت مدرس}
\faculty{مهندسی برق و کامپیوتر}
\department{گروه هوش مصنوعی و رباتیکز}
\subject{مهندسی کامپیوتر}
\field{نرم‌افزار}
\credits{---} % در رساله دکتری معمولا به صورت واحدی ذکر نمی‌شود یا فیلد خالی است
\title{یادگیری تقویتی عمیق با شکل‌دهی تصادفی پاداش برای مدیریت عدم قطعیت در سیستم‌های خودوفقی}
\firstsupervisor{دکتر ...}
% \secondsupervisor{استاد راهنمای دوم}
\firstadvisor{دکتر ...}
% \secondadvisor{استاد مشاور دوم}
\name{علیرضا}
\surname{کاویانی‌فر}
\studentID{...}
\thesisdate{بهمن ۱۴۰۴}

% مشخصات انگلیسی
\latinuniversity{Tarbiat Modares University}
\latinfaculty{Faculty of Electrical and Computer Engineering}
\latinsubject{Computer Engineering}
\latinfield{Software}
\latintitle{Randomized Reward-Shaped Deep Reinforcement Learning for Managing Uncertainty in Self-Adaptive Systems}
\firstlatinsupervisor{Dr. ...}
\firstlatinadvisor{Dr. ...}
\latinname{Alireza}
\latinsurname{Kavianifar}
\latinthesisdate{February 2026}
\latinkeywords{Self-Adaptive Systems, Deep Reinforcement Learning, Reward Shaping, Uncertainty}

% فراخوانی بسته‌ها و دستورات استاندارد
% در این فایل، دستورها و تنظیمات مورد نیاز، آورده شده است.
%-------------------------------------------------------------------------------------------------------------------

% در ورژن جدید زی‌پرشین برای تایپ متن‌های ریاضی، این سه بسته، حتماً باید فراخوانی شود
\usepackage{amsthm,amssymb,amsmath}
% بسته‌ای برای تنطیم حاشیه‌های بالا، پایین، چپ و راست صفحه
\usepackage[top=30mm, bottom=25mm, left=25mm, right=35mm]{geometry}
% بسته‌‌ای برای ظاهر شدن شکل‌ها و تصاویر متن
\usepackage{graphicx}
% بسته‌ای برای رسم کادر
\usepackage{framed} 
% بسته‌‌ای برای چاپ شدن خودکار تعداد صفحات در صفحه «معرفی پایان‌نامه»
\usepackage{lastpage}
% بسته‌ و دستوراتی برای ایجاد لینک‌های رنگی با امکان جهش
% \usepackage[pagebackref=false,colorlinks,linkcolor=black,citecolor=black]{hyperref}
% چنانچه قصد پرینت گرفتن نوشته خود را دارید، خط بالا را غیرفعال و  از دستور زیر استفاده کنید چون در صورت استفاده از دستور زیر‌‌، 
% لینک‌ها به رنگ سیاه ظاهر خواهند شد که برای پرینت گرفتن، مناسب‌تر است
% \usepackage[pagebackref=false]{hyperref}
% بسته‌ لازم برای تنظیم سربرگ‌ها
\usepackage{fancyhdr}
\setlength{\headheight}{16pt}
%
\usepackage{setspace}
\usepackage{algorithm}
\usepackage{algorithmic}
\usepackage{subfigure}
\usepackage{titletoc}


% بسته‌ای برای ظاهر شدن «مراجع» و «نمایه» در فهرست مطالب
\usepackage[nottoc]{tocbibind}
% دستورات مربوط به ایجاد نمایه
\usepackage{makeidx}
\makeindex
% تعریف شیم‌های مورد نیاز برای نسخه‌های جدید بسته bidi در برخی محیط‌ها
% Removal of redundant and potentially broken definitions
\makeatletter
% شیم‌های سازگاری برای نسخه‌های جدید لاتک و بسته‌های bidi/xepersian
\providecommand{\IfClassLoadedT}[2]{\@ifclassloaded{#1}{#2}{}}
\providecommand{\IfClassLoadedF}[2]{\@ifclassloaded{#1}{}{#2}}
\providecommand{\IfClassLoadedTF}[3]{\@ifclassloaded{#1}{#2}{#3}}
\providecommand{\IfPackageLoadedT}[2]{\@ifpackageloaded{#1}{#2}{}}
\providecommand{\IfPackageLoadedF}[2]{\@ifpackageloaded{#1}{}{#2}}
\providecommand{\IfPackageLoadedTF}[3]{\@ifpackageloaded{#1}{#2}{#3}}
\providecommand{\@iffileloaded}[3]{\expandafter\ifx\csname ver@#1\endcsname\relax#3\else#2\fi}
\providecommand{\IfFileLoadedT}[2]{\@iffileloaded{#1}{#2}{}}
\providecommand{\IfFileLoadedF}[2]{\@iffileloaded{#1}{}{#2}}
\providecommand{\IfFileLoadedTF}[3]{\@iffileloaded{#1}{#2}{#3}}
\makeatother

%%%%%%%%%%%%%%%%%%%%%%%%%%
% فراخوانی بسته زی‌پرشین و تعریف قلم فارسی و انگلیسی
\usepackage[pagebackref=false,colorlinks,linkcolor=black,citecolor=black]{hyperref}
\usepackage{notoccite}

\ExplSyntaxOn
\cs_if_exist:NTF \g__hyp_linknestlevel_int { } { \int_new:N \g__hyp_linknestlevel_int }
\cs_if_exist:NTF \l__hyp_annot_GoTo_bool { } { \bool_new:N \l__hyp_annot_GoTo_bool }
\cs_if_exist:NTF \l__hyp_annot_URI_bool { } { \bool_new:N \l__hyp_annot_URI_bool }
\cs_if_exist:NTF \l__hyp_href_url_ismap_bool { } { \bool_new:N \l__hyp_href_url_ismap_bool }
\cs_if_exist:NTF \l__hyp_text_enc_uri_print_tl { } { \tl_new:N \l__hyp_text_enc_uri_print_tl }
\cs_if_exist:NTF \l__hyp_uri_tmpa_tl { } { \tl_new:N \l__hyp_uri_tmpa_tl }

% Defining missing l3pdf/pdfdict commands if not present
\cs_if_exist:NTF \pdfdict_put:nno { } { \cs_new:Npn \pdfdict_put:nno #1#2#3 { } }
\cs_if_exist:NTF \pdfdict_put:nnn { } { \cs_new:Npn \pdfdict_put:nnn #1#2#3 { } }
\cs_if_exist:NTF \pdfdict_use:n { } { \cs_new:Npn \pdfdict_use:n #1 { } }
\cs_if_exist:NTF \pdfannot_dict_put:nne { } { \cs_new:Npn \pdfannot_dict_put:nne #1#2#3 { } }
\cs_if_exist:NTF \pdfannot_link:nen { } { \cs_new:Npn \pdfannot_link:nen #1#2#3 { } }
\cs_if_exist:NTF \__hyp_check_link_nesting:TF { } { \cs_new:Npn \__hyp_check_link_nesting:TF #1#2 { #1 } }
\cs_if_exist:NTF \__hyp_text_pdfstring:eoN { } { \cs_new:Npn \__hyp_text_pdfstring:eoN #1#2#3 { } }
\ExplSyntaxOff

% تعریف دستورات گمشده برای جلوگیری از خطای renewcommand در بسته‌های داخلی
\makeatletter
\providecommand{\foottextfont}{}
\providecommand{\LTRfoottextfont}{}
\providecommand{\RTLfoottextfont}{}
\providecommand{\abstractname}{}
\providecommand{\latinabstract}{}
\providecommand{\AutoPersianDigits}{}
\makeatother

\usepackage{xepersian}
\settextfont[Path=font/,Extension=.ttf,Scale=1,
    BoldFont=XB NiloofarBd,
    ItalicFont=XB NiloofarIt,
    BoldItalicFont=XB NiloofarBdIt
]{XB Niloofar}
\setlatintextfont[Scale=0.9]{Times New Roman}

%%%%%%%%%%%%%%%%%%%%%%%%%%
% چنانچه می‌خواهید اعداد در فرمول‌ها، انگلیسی باشد، خط زیر را غیرفعال کنید
\setdigitfont[Path=font/,Extension=.ttf,Scale=1,
    BoldFont=XB NiloofarBd,
    ItalicFont=XB NiloofarIt,
    BoldItalicFont=XB NiloofarBdIt
]{XB Niloofar}
%%%%%%%%%%%%%%%%%%%%%%%%%%
% تعریف قلم‌های فارسی و انگلیسی اضافی برای استفاده در بعضی از قسمت‌های متن
\defpersianfont\titlefont[Path=font/,Extension=.ttf,Scale=1]{XB Titre}
% \defpersianfont\iranic[Scale=1.1]{XB Zar Oblique}%Italic}%
% \defpersianfont\nastaliq[Scale=1.2]{IranNastaliq}

%%%%%%%%%%%%%%%%%%%%%%%%%%
% دستوری برای حذف کلمه «چکیده»
\renewcommand{\abstractname}{}
% دستوری برای حذف کلمه «abstract»
%\renewcommand{\latinabstract}{}
% دستوری برای تغییر نام کلمه «اثبات» به «برهان»
\renewcommand\proofname{\textbf{برهان}}
% دستوری برای تغییر نام کلمه «کتاب‌نامه» به «مراجع»
\renewcommand{\bibname}{مراجع}
% دستوری برای تعریف واژه‌نامه انگلیسی به فارسی
\newcommand\persiangloss[2]{#1\dotfill\lr{#2}\\}
% دستوری برای تعریف واژه‌نامه فارسی به انگلیسی 
\newcommand\englishgloss[2]{#2\dotfill\lr{#1}\\}
% تعریف دستور جدید «\پ» برای خلاصه‌نویسی جهت نوشتن عبارت «پروژه/پایان‌نامه/رساله»
\newcommand{\پ}{پروژه/پایان‌نامه/رساله }

%\newcommand\BackSlash{\char`\\}

%%%%%%%%%%%%%%%%%%%%%%%%%%
\SepMark{-}

% تعریف و نحوه ظاهر شدن عنوان قضیه‌ها، تعریف‌ها، مثال‌ها و ...
\theoremstyle{definition}
\newtheorem{definition}{تعریف}[section]
\theoremstyle{plain}
\newtheorem{theorem}[definition]{قضیه}
\newtheorem{lemma}[definition]{لم}
\newtheorem{proposition}[definition]{گزاره}
\newtheorem{corollary}[definition]{نتیجه}
\newtheorem{remark}[definition]{ملاحظه}
\theoremstyle{definition}
\newtheorem{example}[definition]{مثال}

%\renewcommand{\theequation}{\thechapter-\arabic{equation}}
%\def\bibname{مراجع}
\numberwithin{algorithm}{chapter}
\def\listalgorithmname{فهرست الگوریتم‌ها}
\def\listfigurename{فهرست تصاویر}
\def\listtablename{فهرست جداول}

%%%%%%%%%%%%%%%%%%%%%%%%%%%%
% دستورهایی برای سفارشی کردن سربرگ صفحات
% \newcommand{\SetHeader}{
% \csname@twosidetrue\endcsname
% \pagestyle{fancy}
% \fancyhf{} 
% \fancyhead[OL,EL]{\thepage}
% \fancyhead[OR]{\small\rightmark}
% \fancyhead[ER]{\small\leftmark}
% \renewcommand{\chaptermark}[1]{%
% \markboth{\thechapter-\ #1}{}}
% }
%%%%%%%%%%%%5
%\def\MATtextbaseline{1.5}
%\renewcommand{\baselinestretch}{\MATtextbaseline}
\doublespacing
%%%%%%%%%%%%%%%%%%%%%%%%%%%%%
% دستوراتی برای اضافه کردن کلمه «فصل» در فهرست مطالب با استفاده از titletoc
\titlecontents{chapter}
  [5em]
  {\addvspace{1em}\bfseries}
  {\contentslabel[\chaptername~\thecontentslabel:]{5em}}
  {\hspace*{-5em}}
  {\titlerule*[0.5pc]{.}\contentspage}

\titlecontents{section}
  [5em]
  {}
  {\contentslabel{2em}}
  {\hspace*{-2em}}
  {\titlerule*[0.5pc]{.}\contentspage}

\titlecontents{subsection}
  [8em]
  {}
  {\contentslabel{3em}}
  {\hspace*{-3em}}
  {\titlerule*[0.5pc]{.}\contentspage}

\titlecontents{subsubsection}
  [12em]
  {}
  {\contentslabel{4em}}
  {\hspace*{-4em}}
  {\titlerule*[0.5pc]{.}\contentspage}

\makeatletter
{\def\thebibliography#1{\chapter*{\refname\@mkboth
   {\uppercase{\refname}}{\uppercase{\refname}}}\list
   {[\arabic{enumi}]}{\settowidth\labelwidth{[#1]}
   \rightmargin\labelwidth
   \advance\rightmargin\labelsep
   \advance\rightmargin\bibindent
   \itemindent -\bibindent
   \listparindent \itemindent
   \parsep \z@
   \usecounter{enumi}}
   \def\newblock{}
   \sloppy
   \sfcode`\.=1000\relax}}
\makeatother


% تنظیمات اختصاصی برای فهرست مطالب (TOC)
\setcounter{secnumdepth}{3}
\setcounter{tocdepth}{3}

% بازگرداندن نقاط رابط (Dot Leaders) که در Cleanup حذف شده بود
\renewcommand{\cftchapleader}{\cftdotfill{\cftdotsep}}
\renewcommand{\cftsecleader}{\cftdotfill{\cftdotsep}}
\renewcommand{\cftsubsecleader}{\cftdotfill{\cftdotsep}}

% تنظیمات تورفتگی و ظاهر
\setlength{\cftsubsubsecindent}{6em}
\setlength{\cftaftertoctitleskip}{10pt} % رفع فاصله زیاد بعد از عنوان فهرست
\raggedbottom

\begin{document}
\AutoPersianDigits
\setstretch{1.5}

\pagenumbering{harfi}
\firstPage
\besmPage
\davaranPage
% در این قسمت اسامی داوران در فایل اصلی باید دستی وارد شوند، اینجا از ساختار پیش‌فرض استفاده می‌کنیم

\esalatPage
\mojavezPage

\chapter*{تقدیم}
\addcontentsline{toc}{chapter}{تقدیم}
\begin{center}
    \textit{(متن تقدیم به خانواده، اساتید و ...)}
\end{center}
\newpage

\begin{acknowledgementpage}
(متن تشکر و قدردانی از اساتید راهنما، مشاور و دوستان)
\signature 
\end{acknowledgementpage}

\keywords{سیستم‌های خودوفقی، یادگیری تقویتی عمیق، شکل‌دهی پاداش، عدم قطعیت}
\fa-abstract{
(خلاصه‌ای از چالش، راهکار و نتایج کلیدی در ۳۰۰ تا ۵۰۰ کلمه)
}
\abstractPage

\tableofcontents
\newpage

\listoftables
\newpage

\listoffigures
\newpage

\chapter*{فهرست علائم و اختصارات}
\addcontentsline{toc}{chapter}{فهرست علائم و اختصارات}
(لیست نمادهای ریاضی مانند $\rho$، $\gamma$ و اختصارات مانند \lr{DRL}، \lr{SAS})
\newpage

% ============================================================================
% متن اصلی رساله
% ============================================================================
\pagenumbering{arabic}
\pagestyle{fancy}

% فصل ۱: مقدمه
\chapter{مقدمه}
\section{مقدمه}

در حوزه مهندسی سیستم‌های نرم‌افزاری معاصر، پیچیدگی به عنوان یکی از بارزترین ویژگی‌ها خودنمایی می‌کند. محیط عملیاتی پویا نقش محوری در شکل‌دهی به این پیچیدگی ایفا کرده و سیستم‌ها را مجبور می‌سازد تا با شرایط نامشخصی که معمولاً تنها پس از عملیاتی شدن بروز می‌کنند، دست و پنجه نرم کنند.

علاوه بر این، چالش‌ها و منابع عدم قطعیت در سیستم‌های مختلف به شکل‌های گوناگونی نمود پیدا می‌کنند. به عنوان مثال، در سیستم‌های مستقر در محیط‌های پویا نظیر برنامه‌های اینترنت اشیا (\lr{IoT}) در شهرهای هوشمند، تغییرات آب و هوایی تأثیر قابل توجهی می‌گذارد. همچنین محیط‌های ابری با تقاضای منابع متغیر، امنیت سایبری با تهدیدات نوظهور و مداوم، تجارت الکترونیک و شبکه‌های اجتماعی با رفتار غیرقابل پیش‌بینی کاربران، و سیستم‌های توزیع‌شده با تأخیر شبکه، خرابی سخت‌افزار و تغییرات زیرساختی روبه‌رو هستند.

سیستم‌های سنتی ایستا یا مبتنی بر قواعد، انعطاف‌پذیری لازم برای پاسخ به چنین تغییراتی در زمان اجرا را ندارند. در مقابل، سیستم‌های خودوفقی (\lr{SASs}) مجهز به حلقه‌های بازخوردی هستند که به آن‌ها اجازه می‌دهد تا محیط و وضعیت خود سیستم را پایش کنند، انحرافات از رفتار مورد انتظار را تحلیل نمایند و به طور خودمختار استراتژی‌ها و تغییرات وفقی را برنامه‌ریزی و اجرا کنند.

هنگامی که سیستم با یک فضای وفقی و گزینه‌های تغییر بسیار گسترده مواجه می‌شود، وظیفه انتخاب یک گزینه وفقی که با اهداف سیستم همسو باشد، از نظر محاسباتی بسیار پرهزینه و دشوار، و در برخی موارد تقریباً غیرممکن می‌گردد. تمرکز این رساله بر سیستم‌هایی است که دارای فضای وفقی گسسته و بزرگ هستند و نیاز فوری به راهکارهای کارآمد جهت انتخاب گزینه‌های بهینه در زمان اجرا دارند.

\section{بیان مسئله و سوالات تحقیق}

\subsection{مسئله تحقیق}

روش‌های موجود برای مدیریت عدم قطعیت در سیستم‌های خودوفقی با محدودیت‌های قابل توجهی روبه‌رو هستند. به عنوان مثال، روش‌های مبتنی بر یادگیری ماشین نیازمند داده‌های آموزشی برچسب‌دار هستند که نمایانگر محیط سیستم باشند؛ به دست آوردن چنین داده‌هایی در زمان طراحی به دلیل عدم قطعیت فراوان بسیار دشوار است. از طرفی، روش‌های مبتنی بر جستجو (مانند الگوریتم‌های صعود تپه و الگوریتم‌های ژنتیک) با مشکلات مقیاس‌پذیری مواجه هستند و عملکرد آن‌ها با تغییر شرایط محیطی دچار افت محسوس می‌شود.

یادگیری تقویتی عمیق (\lr{DRL}) اگرچه ابزاری قدرتمند است، اما خود با چالش‌های اساسی مانند گیر افتادن در بهینه‌های محلی مواجه می‌باشد. بهینه‌های محلی به راه‌حل‌هایی اشاره دارند که تنها در یک منطقه محدود از فضای مسئله بهینه هستند و به صورت سراسری بهترین راه‌حل ممکن محسوب نمی‌شوند. در زمینه \lr{DRL}، این بدان معناست که مدل ممکن است سیاستی را یاد بگیرد که تحت شرایط خاص به خوبی عملکرد مناسبی دارد، اما در صورت تغییر محیط قادر به تعمیم‌پذیری و حفظ عملکرد مطلوب نیست.

علاوه بر این، آموزش مجدد یا بازآموزی (\lr{Retraining}) مدل در زمان اجرا فرآیندی پرهزینه است. بازآموزی شامل مراحل متعددی مانند جمع‌آوری داده‌های جدید، ارزیابی مجدد ساختار و پارامترهای مدل، و تست مدل به‌روزشده برای اطمینان از عملکرد بهتر آن نسبت به نسخه قبلی است. در طول این فرآیند، سیستم ممکن است به صورت غیربهینه عمل کند که در سیستم‌های حیاتی، این دوره عملکرد نامطلوب می‌تواند پیامدهای جبران‌ناپذیری به همراه داشته باشد.

\textbf{بیان اصلی و خلاصه‌شده‌ی مسئله:} ماهیت پویا و غیرقابل پیش‌بینی محیط‌هایی که سیستم‌های نرم‌افزاری در آن‌ها عمل می‌کنند، چالش‌های قابل توجهی را برای مدل‌های سنتی یادگیری تقویتی عمیق به وجود می‌آورد. این مدل‌ها به طور مکرر با مشکل بهینه‌های محلی دست و پنجه نرم کرده و نمی‌توانند در زمان مواجهه با شرایط متغیر به خوبی عملکرد خود را تعمیم دهند که این امر منجر به عملکرد نامطلوب سیستم و نیاز به بازآموزی‌های مکرر و زمان‌بر خواهد شد.

\subsection{سوالات تحقیق}

با توجه به چالش‌ها و خلاءهای تحقیقاتی مطرح شده و بر اساس دستاوردهای مقاله پیشین \cite{RS_DRL_Article}، این رساله به دنبال پاسخگویی به سوالات اصلی زیر است:

\begin{enumerate}
    \item \textbf{سوال اول:} چگونه می‌توان یک روش مبتنی بر یادگیری تقویتی عمیق طراحی کرد که به طور موثری از بهینه‌های محلی در محیط‌های پویا، بدون نیاز به بازآموزی مکرر، فرار کند؟
    \item \textbf{سوال دوم:} چگونه می‌توان مکانیزم‌های شکل‌دهی پاداش (\lr{Reward Shaping}) را با بازپخش تجربه (\lr{Experience Replay}) ادغام کرد تا تعمیم‌پذیری سیستم در تغییرات شرایط محیطی بهبود یابد؟
    \item \textbf{سوال سوم:} چگونه روش پیشنهادی می‌تواند فضاهای وفقی گسسته و بزرگ (از ۲۱۶ تا ۴۰۹۶ گزینه) را در حین حفظ عملکرد نزدیک به بهینه، به شکل کارآمدی مدیریت کند؟
    \item \textbf{سوال چهارم:} تا چه حد روش پیشنهادی (\lr{RS-DRL}) از روش‌های موجود نظیر یادگیری ای‌گرایانه حریصانه (\lr{$\epsilon$-greedy DQN})، \lr{DLASER} و \lr{ML4EAS} از منظر عملکرد مجانبی، سرعت همگرایی، و بهینه‌سازی چندهدفه، برتری و عملکرد بهتری از خود نشان می‌دهد؟
\end{enumerate}

\section{اهداف رساله}

این رساله اهداف کلان و خرد زیر را در راستای توسعه تحقیقات پیشین \cite{RS_DRL_Article} دنبال می‌کند:

\begin{enumerate}
    \item \textbf{طراحی و توسعه:} طراحی یک معماری نوآورانه یادگیری تقویتی عمیق با شکل‌دهی پاداش تصادفی (\lr{RS-DRL}) که به شکل یکپارچه‌ای با حلقه خودمختار \lr{MAPE-K} در سیستم‌های خودوفقی ادغام شده باشد.
    \item \textbf{نوآوری الگوریتمی:} توسعه یک مکانیزم نوین شکل‌دهی پاداش که تجربیات ناموفق را در طول فرآیند بازپخش حافظه به صورت تصادفی بازطراحی می‌نماید تا اکتشاف (\lr{Exploration}) را ترویج کرده و از همگرایی مدل در بهینه‌های محلی جلوگیری کند.
    \item \textbf{مدل‌سازی صوری:} مدل‌سازی صوری مسئله کنترل و وفقی‌سازی تحت عدم قطعیت که شامل فضای حالت (نظیر سطح انرژی، از دست رفتن بسته‌های شبکه، زمان تأخیر)، فضای عمل (گزینه‌های وفقی گسسته) و توابع پاداش جهت بهینه‌سازی‌های تک‌هدفه و چندهدفه است.
    \item \textbf{ارزیابی تجربی:} ارزیابی راهکار توسعه‌یافته \lr{RS-DRL} در تقابل با روش‌های پیشرفته‌ی فعلی (مانند \lr{Soft HER}، \lr{HER}، \lr{$\epsilon$-greedy DQN}، \lr{DLASER} و \lr{ML4EAS}) بر روی مطالعات موردی سیستم \lr{DeltaIoT} با اندازه‌های مختلف از فضاهای تطبیق (از جمله ۲۱۶، ۱۲۹۶ و ۴۰۹۶ گزینه ممکن).
    \item \textbf{تحلیل‌های آماری:} ارائه ارزیابی‌های آماری دقیق (تست مستقل \lr{t-test}، مقادیر \lr{p-value} و فاصله‌های اطمینان) برای تصدیق برتری روش پیشنهادی.
\end{enumerate}

\section{نوآوری‌های رساله}

\subsection{نوآوری اول: ارائه الگوریتم یادگیری تقویتی عمیق با شکل‌دهی تصادفی پاداش (\lr{RS-DRL})}
برخلاف بازپخش تجربه سنتی در یادگیری تقویتی، الگوریتم پیشنهاد شده در این رساله (\lr{RS-DRL}) برخی از تجربیات ناموفق گذشته را بازطراحی کرده و با اعطای پاداش مثبت به آن‌ها (همانند پاداش‌های موفقیت)، فرآیند یادگیری انسانی از اشتباهات گذشته خود را شبیه‌سازی می‌کند. مکانیزم انتخاب تصادفی تجربیات ناموفق که به وسیله فاکتور شکل‌دهی $\rho$ کنترل می‌شود، تجربیات مردود سیستم را با تخصیص پاداش خوش‌بینانه (به عنوان مثال تخصیص $r_i \leftarrow 1$) شبیه‌سازی موفقیت می‌نماید. این ترفند برای محیط‌های متنوع نیازی به تعاریف تخصصی و دانش دامنه خاص ندارد و موجب می‌شود تا: ۱) تعمیم‌پذیری را در طول محیط‌های متنوع تشویق نماید و ۲) فرآیند یادگیری را از طریق تبدیل اشتباهات گذشته به درس‌های با ارزش در تجربه‌های بعدی، به شدت تسریع کند.
در مقایسه با روش‌هایی نظیر بازپخش تجربه با نگاه به گذشته (\lr{HER}) و \lr{Soft HER} که ملزم به داشتن دانش دامنه‌ای و داشتن هدف صریح برای برچسب‌زنی مجدد اهداف (\lr{Goal Relabeling}) هستند، رویکرد ما در \lr{RS-DRL} از استراتژی پاداش تصادفی و خوش‌بینانه و کاملاً مستقل از دانش حوزه‌ای بهره می‌برد.

\subsection{نوآوری دوم: یکپارچه‌سازی \lr{RS-DRL} با معماری خودوفقی \lr{MAPE-K}}
نوآوری استراتژیک دیگر، ارائه چارچوب معماری تلفیقی شامل سیستم مدیریتی و سیستم مدیریت‌شونده است. معماری پیشنهادشده که با حلقه کنترل و بازخورد \lr{MAPE-K} ادغام شده است، شامل چرخه ارزیابی و انتخاب ۹ مرحله‌ای در تعامل بین مؤلفه‌های مختلف آن در زمان اجرا خواهد بود.
این اجزای کلیدی شامل ماژول استخراج‌کننده ویژگی از داده‌های سنسورها، انتخابگر فضای گزینه (\lr{$\epsilon$-greedy DQN})، تایید‌کننده گزینه وفقی از طریق اعتبارسنجی مدل با موتور \lr{UPPAAL-SMC}، محاسبه‌گر پاداش، حافظه بازپخش با در نظر گرفتن مکانیزم بدیع شکل‌دهی پاداش و مؤلفه آموزش مدل می‌باشد.

\subsection{نوآوری سوم: مدیریت موثر عدم قطعیت در فضاهای وفقی بزرگ و گسسته}
به عنوان یک نوآوری اساسی در سطح ارزیابی و دستاوردهای سیستم، این راهکار نشان داده که پتانسیل چشم‌گیری در مدیریت و راهبری سیستم‌های با مقیاس‌پذیری فضاهای حالت و فضای عمل بالا، مانند شبکه‌های گسترده‌ی اینترنت اشیاء داراست.
کاربست رویکرد \lr{RS-DRL} در پلتفرم \lr{DeltaIoT} از طریق کاهش چشم‌گیر هدررفت داده‌ها تا محدوده ۲۲.۲۳٪ نسبت به رقبای سرسختی مانند \lr{DLASER+} و بهبود ۳۹.۰۷ درصدی زمان تأخیر انتقال در شبکه، قابل توجه است. استفاده از فضای گسسته وفقی مقیاس‌پذیر در محیط‌های آزمایشگاهی شامل نسخه‌ی \lr{DeltaIoTv1} با ۲۱۶ گزینه، نسخه‌ی \lr{DeltaIoTv1.1} شامل ۱۲۹۶ گزینه و نسخه‌ی \lr{DeltaIoTv2} شامل ۴۰۹۶ حالت ممکن، کارایی و میزان پاداش مجانبی بهبود یافته و پایداری را نسبت به روش‌های پیشین نشان می‌دهد.

\section{مروری بر فصول رساله}

ساختار رساله به گونه‌ای طراحی شده است که به طور نظام‌مند روند توسعه، اجرا و ارزیابی رویکرد پیشنهادی را در اختیار خواننده قرار دهد. در ادامه، مرور خلاصه‌ای از هر یک از فصول آورده شده است:

\textbf{فصل ۲: مبانی و پیشینه تحقیق}\\
مفاهیم بنیادی مربوط به سیستم‌های خودوفقی، چالش‌های عدم قطعیت و مبانی مدل‌سازی آن در زمان اجرا، اصول سیستم‌های یادگیری تقویتی (\lr{MDP}، یادگیری \lr{Q}) و \lr{DRL} (به ویژه الگوریتم‌های \lr{DQN}) به تفصیل شرح داده می‌شود. سپس کارهای مرتبط پیشین در مدیریت فضاهای بزرگ تطبیق، از روش‌های مبتنی بر یادگیری ماشین نظیر \lr{ML4EAS} گرفته تا الگوریتم‌های اکتشافی جستجو و سایر روش‌های نوین مورد بررسی و ارزیابی قرار می‌گیرند.

\textbf{فصل ۳: مدل سیستم و فرمول‌بندی مسئله}\\
در این فصل، سیستم شبکه‌ی سنسورهای اینترنت اشیاء (\lr{DeltaIoT}) به عنوان مطالعه‌ی موردی انگیزشی معرفی می‌شود. فضای حالت صوری شبکه شامل مفاهیمی پیرامون مصرف انرژی، از دست‌رفتگی بسته‌های داده و زمان تأخیر ($S = \{\text{\lr{Energy}}, \text{\lr{Packet Loss}}, \text{\lr{Latency}}\}$) فرمول‌بندی شده و در ادامه، گزینه‌های فضای اعمال گسسته نیز مشخص و تدوین می‌گردد. در نهایت، توابع پاداش مربوط به سناریوهای بهینه‌سازی چندهدفه و تک‌هدفه به صورت ریاضی ارائه می‌شوند.

\textbf{فصل ۴: روش پیشنهادی \lr{RS-DRL}: معماری و الگوریتم}\\
مرکز ثقل این رساله در این فصل واقع شده است که در آن، معماری پیشنهادی یکپارچه با کنترل‌کننده‌ی \lr{MAPE-K} تبیین می‌گردد. الگوریتم پایه (\lr{RS-DRL})، به انضمام جزئیات دقیق مکانیزم شکل‌دهی پاداش‌های تصادفی با جزئیات الگوریتمی تشریح گردیده و پیاده‌سازی‌های مرتبط با اعتبارسنجی استراتژی‌های اجرایی پیش از پیاده‌سازی نهاییِ تغییرات توسط نرم‌افزار \lr{UPPAAL-SMC}، برای تضمین پایداری سیستم، بحث و بررسی می‌شوند.

\textbf{فصل ۵: ارزیابی تجربی}\\
نحوه تخصیص محیط شبیه‌سازی، سناریوهای متعدد پیکربندی شبکه، و پروتکل‌های مقایسه‌ی جامع مدل \lr{RS-DRL} با استراتژی‌های به‌روزِ رقیب نظیر \lr{$\epsilon$-Greedy DQN}، \lr{HER}، \lr{Soft HER}، \lr{DLASER}، و \lr{ML4EAS} در این فصل ارائه می‌گردد. ارائه‌ی شفافِ معیارهای کلیدی دستاوردها شامل بررسی میزان ظرفیت عملکرد مجانبی سیستم، ارزیابی زمان مورد نیاز تقربی برای رسیدن به آستانه قابل اتکا، و در نهایت کل مقادیر پاداش دریافتی در هر دوره‌ی آزمایش همراه با تأییدیه‌های تحلیل آماری و علمی مستند می‌باشند.

\textbf{فصل ۶: بحث، کارهای آینده و نتیجه‌گیری}\\
استلزامات نظریِ مرتبط با علت تفوقِ راهکار پیشنهادی بازبینی گردیده و کارایی مکانیزم از زوایای پراتیک صنعتی بحث می‌شود. در ادامه‌ی این فصل، به بررسی مصالحه‌ها، محدودیت‌ها و چالش‌های باقی‌مانده‌ی این معماری پرداخته شده و در نهایت، مسیرها و چارچوب‌هایی الهام‌بخش جهت بسط رویکرد فعلی در سیستم‌های پیوسته‌ی سایبرفیزیکی آینده پیشنهاد می‌گردد. فصل با نتیجه‌گیری کلی پیرامون مزیت‌های کاربردی راهکار نوین برای سیستم‌های عملیاتی بسته خواهد شد.

\newpage

% فصل ۲: مبانی و پیشینه تحقیق
\chapter{مبانی و پیشینه تحقیق}
\section{مقدمه}

هدف این فصل، ایجاد و بسط یک چارچوب و زیربنای نظری مستحکم است تا خواننده را از مفاهیم سطح بالا به سمت درک دقیق شکاف تحقیقاتی که این رساله آن را پر می‌کند، هدایت نماید. سیستم‌های نرم‌افزاری مدرن همواره با عدم قطعیت‌های مداوم مواجه هستند و باید در زمان اجرا به طور پیوسته ساختار و رفتار خود را با شرایط در حال تغییر وفق دهند. این فصل در ابتدا اصول اولیه سیستم‌های خودوفقی و مؤلفه‌های آن را تشریح می‌کند، سپس به نقش یادگیری تقویتی و محدودیت‌های روش‌های سنتی در مقابله با فضاهای بزرگ تصمیم‌گیری می‌پردازد و در نهایت با مروری ساختاریافته بر ادبیات تحقیق، علت نیاز به راهکار نوین \lr{RS-DRL} را اثبات می‌نماید.

\section{مبانی سیستم‌های خودوفقی}
\subsection{مفاهیم اصلی و الگوهای معماری}

سیستم‌های خودوفقی (\lr{Self-Adaptive Systems}) یک رویکرد برجسته و مؤثر برای مقابله با عدم قطعیت به شمار می‌روند. ایده کلیدی در این سیستم‌ها، تجهیز آن‌ها به یک یا چند حلقه بازخوردی کنترلی است که به صورت مداوم مؤلفه‌های داخلی سیستم و محیط خارجی آن را پایش کرده و در صورت نیاز سیستم را با تغییرات منطبق می‌نمایند \cite{RS_DRL_Article}. 

مهمترین و شناخته‌شده‌ترین الگوی معماری برای پیاده‌سازی این سیستم‌ها، حلقه خودمختار \lr{MAPE-K} است. همان‌طور که در فصل‌های بعدی با جزئیات بیشتر بررسی خواهد شد، این حلقه در کنار سیستم هدف و منطق کسب‌وکار، جریان کنترلی زیر را اعمال می‌کند:
\begin{itemize}
    \item \textbf{پایش (\lr{Monitor}):} سنسورها وضعیت محیط فیزیکی، زیرساخت و خود سیستم را به صورت زمان‌حقیقی استخراج می‌کنند.
    \item \textbf{تحلیل (\lr{Analyze}):} داده‌های پایش‌شده را پردازش کرده و مشخص می‌کند که آیا انحرافی از اهداف کیفی و کمی سیستم رخ داده است و نیاز به اعمال تغییر وجود دارد یا خیر.
    \item \textbf{برنامه‌ریزی (\lr{Plan}):} در صورت نیاز تطبیق، راهکارها و استراتژی‌های موجود (گزینه‌های وفقی) مقایسه شده و گزینه بهینه برای شرایط فعلی برگزیده می‌شود.
    \item \textbf{اجرا (\lr{Execute}):} استراتژی برنامه‌ریزی‌شده در قالب یک سری از دستورات اجرایی از طریق عملگرها (Actuators) بر روی سیستم اعمال می‌گردد.
    \item \textbf{دانش (\lr{Knowledge}):} به عنوان حافظه مرکزی عمل کرده و اطلاعات، مدل‌ها و داده‌های بین چهار مرحله قبلی را هماهنگ و نگهداری می‌کند. راهکار پیشنهادی ما (\lr{RS-DRL}) مستقیماً به عنوان هسته اصلی بخش برنامه‌ریزی در این الگو فعالیت می‌کند.
\end{itemize}

\subsection{چالش‌های عدم قطعیت در زمان اجرا}

عدم قطعیت در سیستم‌های خودوفقی، به هرگونه انحراف از دانش قطعی اطلاق می‌شود که می‌تواند باعث کاهش اطمینان نسبت به تصمیمات زمان اجرا گردد \cite{RS_DRL_Article}. 

به عنوان نمونه‌های عینی، می‌توان به موارد گوناگونی در دامنه‌های مختلف اشاره کرد: در برنامه‌های مبتنی بر اینترنت اشیاء (\lr{IoT}) که در محیط‌های شهری پراکنده هستند، تغییرات مداوم و غیرقابل پیش‌بینی آب و هوا منجر به افت کیفیت سیگنال و اختلال در ارسال اطلاعات سنسورها می‌شود. در محیط‌های ابری، نوسان شدید در تقاضای منابع از سوی کاربران، و در حوزه‌های امنیتی، ظهور بی‌وقفه و تصادفی حملات سایبری از مصادیق برجسته عدم قطعیت هستند.

اهمیت و چالش‌برانگیز بودن عدم قطعیت از آنجا نشأت می‌گیرد که مفروضات، پیکربندی‌ها و تصمیمات در نظر گرفته شده در فاز طراحی معماری نرم‌افزار را بی‌اعتبار می‌سازد. در نتیجه، سیستم باید قادر باشد با یادگیری آنلاین و پویا در زمان اجرا، اهداف اصلی خود را محقق نماید؛ وگرنه با نقض کیفیت خدمات صدمات سنگین احتمالی را تجربه خواهد نمود.

\section{مدل‌سازی عدم قطعیت}
\subsection{منابع عدم قطعیت}

منابع عدم قطعیت به طور کلی در سه دسته وسیع طبقه‌بندی می‌شوند: نخست، نامشخص بودن موجودیت‌های محیطی نظیر پارازیت در حسگرها، نویزها و تغییرات جوی که از کنترل ساختار درونی سیستم خارج هستند. دوم، ماهیت غیرقابل اتکای منابع خود سیستم مثل نقص در قطعات زیربنایی، کاهش سریع سطح باتری یا مشکلات نرم‌افزاری تعاملات میان سرویس‌ها؛ و در نهایت عدم قطعیت در اهداف متغیر و ترجیحات نامشخص ذینفعان در محیط‌های تجاری.

\subsection{رویکردهای مدل‌سازی عدم قطعیت}

جهت درک و محاسبه رفتار سیستم در هنگام بروز مقادیر ناشناخته، استفاده از مدل‌سازی‌های زمان اجرا (\lr{Runtime Models}) ضروری به نظر می‌رسد. در شرایطی که فضای جستجو بسیار بزرگ و ریسک تصمیم‌گیری به دلیل شرایط مبهم بالاست، رویکرد پیشنهادی از مکانیزم‌های مدل‌سازی ریاضی صوری بهره‌گیری و پشتیبانی می‌کند.

در این راستا ابزارهایی مانند موتور \lr{UPPAAL-SMC} و چارچوب \lr{ActivFORMS} به کار گرفته می‌شوند. این زیرساخت‌ها با استفاده از اتوماتای زمان‌دار (\lr{Timed Automata}) مدلی رفتاری از سیستم هدف و محیط آن را توسعه داده و از طریق روش‌های کنترل مدل آماری (\lr{Statistical Model Checking}) می‌توانند به سرعت گزینه‌های وفقی پیشنهاد شده توسط سیستم هوشمند را پیش از اجرا در سیستم فیزیکی ارزیابی، تایید یا رد کنند. این قابلیت نقشی کلیدی در امنیت راهکار پیشنهادی (در بخش تاییدگر گزینه وفقی) این رساله ایفا می‌کند.

\section{یادگیری تقویتی برای وفقی‌سازی}
\subsection{اصول یادگیری تقویتی (\lr{MDP}، یادگیری \lr{Q})}

یادگیری تقویتی از طریق تعامل مستمر یک عامل (\lr{Agent}) با محیط اطرافش بر پایه آزمون و خطا بنا شده است و این محیط به لحاظ ریاضی از طریق فرآیند تصمیم‌گیری مارکوف (\lr{MDP}) فرموله می‌شود. فرآیند \lr{MDP} خود شامل فضای حالت (\lr{State})، فضای عمل (\lr{Action}) و توابع پاداش (\lr{Reward}) است.

در فرآیند تصمیم‌گیری مارکوف، برای هر حالت خاص، عامل عمل مشخصی را اتخاذ می‌کند و مطابق با معادله مشهور بلمن (\lr{Bellman Equation}) محیط ارزش عمل اتخاذ شده را در قالب یک پاداش ارزیابی می‌نماید و به وضعیت بعدی منتقل می‌شود. این مفهوم در سیستم‌های خودوفقی مانند محیط \lr{DeltaIoT} به این صورت اعمال می‌گردد که وضعیت سیستم با معیارهایی نظیر انرژی ذخیره شده سیستم و تاخیر اندازه‌گیری شده فضای حالت را تشکیل می‌دهند و اقدام عامل، همان گزینه وفقی مورد انتخاب خواهد بود و محاسبه پاداش نیز طبق معادله بهینه‌سازی سیستم حاصل می‌گردد. نقش الگوریتم مشهور یادگیری Q ایجاد جدولی جامع برای تخصیص و پیدا کردن ارزش‌هاست تا مدل به سیاست بهینه همگرا شود.

\subsection{یادگیری تقویتی عمیق (\lr{DQN}) برای فضاهای بزرگ}

اما زمانی که با سیستم‌های واقعی که دارای فضای عمل گسسته ولی بسیار پهناور (برای مثال ۴۰۹۶ گزینه وفقی در شبکه‌های اینترنت اشیاء) مواجه هستیم، استفاده از متد کلاسیک یادگیری \lr{Q} با چالش نفرین ابعاد (\lr{Curse of Dimensionality}) روبه‌رو می‌شود؛ زیرا ذخیره و بروزرسانی جدولی با میلیون‌ها حالت و عمل از لحاظ محاسباتی ناکارآمد و در بسیاری موارد غیرممکن است.

شبکه‌های \lr{Q} عمیق (\lr{DQN}) راهکاری هوشمندانه برای این محدودیت ارائه می‌دهند؛ آن‌ها به جای جداول، تابع ارزش کیفیت \lr{Q} را با استفاده از شبکه‌های عصبی عمیق تقریب می‌زنند. برای غلبه بر بی‌ثباتی آموزش شبکه‌های عصبی در \lr{RL}، مدل \lr{DQN} از دو نوآوری برجسته بهره می‌برد: نخست استفاده از مکانیسم حافظه بازپخش تجربه (\lr{Experience Replay}) که نمونه‌های آموزشی تصادفی‌سازی شده و باعث شکستن همبستگی زمانی متوالی می‌گردد؛ و دوم تعبیه هدف ثابت به واسطه شبکه هدف ثابت (\lr{Target Network}) که بهبود تثبیت توابع هزینه را در هنگام محاسبه پیش‌بینی‌ها به دنبال دارد.

\subsection{مفهوم یادگیری مداوم در حلقه‌های وفقی}

اصلی‌ترین کاربرد این رساله، ادغام مفهوم آموزش پیوسته و مداوم توسط عامل \lr{DQN} در چارچوب معماری خودوفقی و به طور خاص درون حلقه \lr{MAPE-K} (همانطور که در فصل ۴ و شکل ۴ از مقاله \cite{RS_DRL_Article} به تفصیل نشان داده خواهد شد) است. پس از جاگذاری در مؤلفه برنامه‌ریزی (\lr{Planner})، عامل یادگیرنده وضعیتی را به عنوان ورودی فضای حالت از ماژول پایش دریافت نموده، استراتژی اقدام را برنامه‌ریزی کرده و بر اساس میزان موفقیت انتخابش از سوی مؤلفه تحلیل پاداش می‌گیرد. این ساختار تعاملی، سیستم خودوفقی را از پیکربندی ایستا به ساختار یادگیری پویا تبدیل می‌سازد.

\section{کارهای مرتبط: مدیریت فضاهای وفقی بزرگ}

دستیابی به کارایی مناسب در مدیریت محیط‌های مبهم بسیار کلیدی است. برای تبیین عمیق این مباحث، در ادامه به مقایسه انتقادی دسته‌بندی‌های موجود می‌پردازیم.

\subsection{روش‌های مبتنی بر یادگیری ماشین (\lr{ML-based Methods})}

اولین دسته‌بندی کارهای مرتبط در به‌کارگیری یادگیری ماشین نظارت‌شده و پیش‌بینی‌کننده با هدف هرس و کاهش فضای جستجو متمرکز است. رویکرد محققانی چون کین (\lr{Quin}) و همکارانش \cite{quin2019efficient,quin2020machine} و مدل یادگیری برای ارزیابی کیفی ماشین (\lr{ML4EAS}) توسعه یافته توسط کاچی و بواناکا \cite{kachi2021machine} نشان می‌دهند که الگوریتم‌های طبقه‌بندی و رگرسیون قادرند تعداد زیادی از گزینه‌های غیربهینه را تا پیش از ارسال به مؤلفه تحلیل‌گر، غربال کنند.

\textbf{نقد رویکرد:} اگرچه این روش‌ها در کاهش اندازه محیط عمل خود کارا هستند، اما وجود وابستگی بنیادین به مجموعه‌داده‌های با برچسب آموزش‌دیده، نقص متدولوژی است. تهیه مجموعه‌داده در زمان توسعه که همسان با متغیرهای غیرمترقبه واقعیت باشد غیرعملی است. متعاقباً با بروز رانش مفهومی (\lr{Concept Drift}) ناشی از تغییر مداوم شرایط محیطی، این مدل‌های پیش‌آموزش دیده، عملکرد تعمیم‌یافته و کارآمد خود را در برخورد با رویه‌های ناشناخته از دست می‌دهند.

\subsection{روش‌های مبتنی بر جستجو (\lr{Search-based Methods})}

متدولوژی کلاسیک دیگر برای پیمایش سریع فضاهای وسیع جستجو استفاده از راهبردهای اکتشافی است. برای مثال استفاده کوکر (\lr{Coker}) از الگوریتم صعود تپه (\lr{Hill Climbing}) \cite{coker2015evaluating} و برنامه‌ریزی‌های مبتنی بر الگوریتم ژنتیک (\lr{Genetic Programming}) که به قلمرو مقالات کینر (\lr{Kinneer}) و همکاران اختصاص دارد \cite{kinneer2018managing,kinneer2018search,kinneer2019reusing}. همچنین شبیه‌سازی جستجو با استفاده از رویکرد تکامل چند هدفه نظیر \lr{MOEA} توسط چن (\lr{Chen}) \cite{chen2018learning}، از کاربردهای متداول است.

\textbf{نقد رویکرد:} هرچند ویژگی انطباق خودکار بدون آموزش در این رویکرد به چشم می‌خورد، اما چالش مقیاس‌پذیری، مانع استقرار کامل آن‌ها در پروژه‌های بلادرنگ است. اجرای مجدد جمعیت تکاملی همچون ژنتیک برای یافتن سناریوی مطلوب در پی هر گونه تطابق زمانی، مستلزم محاسبات بی‌نهایت پرهزینه و ایجاد تاخیر زیاد است؛ موضوعی که در شبکه‌های حساس تاخیرپذیر غیرقابل تحمل است.

\subsection{روش‌های مبتنی بر یادگیری تقویتی (\lr{RL-based Methods})}

پیشرفت‌های حاصل شده در اتوماسیون برخط سیستم‌ها با کارهای کیم و پارک \cite{kim2010reinforcement}، متزگر \cite{metzger2020reinforcement} و مدل \lr{RLMC} توسط کامارا \cite{camara2020reinforcement} به دلیل تعامل بی‌وقفه پویایی‌های آنلاین شناخته شده است که عمدتا براساس کلاسیک الگوریتم \lr{Q-Learning} توسعه یافتند.

\textbf{نقد رویکرد:} الگوریتم کلاسیک در ابعاد بزرگ، همانند شبکه‌ای با میلیاردها ماتریس برای مدل‌سازی از پای در می‌آید و ناتوان از رسیدگی است. حتی تلاش‌هایی که با مدل پایه ارزیابی ما یعنی معماری حریصانه کلاسیک \lr{DQN} با استراتژی $\epsilon$\lr{-greedy} صورت پذیرفته است، کماکان با مشکل شدید محدودیت روبه‌رو است؛ زیرا مدل پایه عامل هوشمندی را پرورش می‌دهد که اغلب در نقاط و قله‌های بهینه‌های محلی (\lr{Local Optima}) متوقف شده، مانع قابلیت یافتن قله‌های اصلی (\lr{Global Optima}) در دینامیک شرایط محیط می‌گردد.

\subsection{شکل‌دهی پاداش و بازپخش تجربه با نگاه به گذشته (\lr{Reward Shaping \& HER})}

بازپخش تجربه‌های ناموفق از موضوعات بسیار داغ در حوزه هوش مصنوعی بوده که مکانیزم مشهور یادگیری با نگاه به گذشته (\lr{HER}) \cite{andrychowicz2017hindsight} و مدل متأخر با رویکرد پاداش نرم‌تر (\lr{Soft HER}) \cite{zheng2023hindsight} با استفاده از استراتژی برچسب‌گذاری مجدد هدف برای تجربیات ناموفق توانسته یادگیری محیط‌ها را سریع‌تر نمایند.

\textbf{نقد رویکرد انطباق با این حوزه:} برتری این مدل‌های عمیق کاملاً محدود به طراحی‌های وابسته به هدف و صریح (\lr{Goal-Conditioned Environments}) می‌باشد. به عبارت دیگر عامل باید یک مکان یا وضعیت شفاف (مانند موقعیت سه‌بعدی ابزار مکانیکی) را تصرف کند؛ در حالی که کنترل‌کننده‌های سیستم‌های خودوفقی اساساً و غالباً اهدافی به عنوان متغیر قطعی را دنبال نمی‌کنند، بلکه نگه‌داشتن انرژی یا تاخیر شبکه در حداقل‌ها یک سیاست آستانه‌محور است که با استراتژی تغییر برچسب و ارزیابی در بازتولید \lr{HER} با سازگاری اساسی تناقض دارد.

\section{سنتز و شناسایی شکاف تحقیقاتی}

با استناد به مطالب شرح داده شده و مروری جامع در کاستی مقالات پیشین، درمی‌یابیم که هیچ‌یک از دسته‌بندی‌های فعلی قادر به برطرف ساختن کل چالش‌ها به صورت جامع نبوده‌اند: راهکارهای یادگیری ماشین به پدیده‌های پیش‌بینی‌نشده محیط به دلیل فقدان آموزش آفلاین به شدت حساس و شکننده‌اند؛ مسیرهای اکتشاف جستجو با تاخیر غیرقابل پذیرش روبرو می‌باشند. حتی اگر از رویکرد‌های \lr{DRL} پایه بهره بگیریم، مستعد توقیف در استثنائات محلی بوده و قابلیت فراگیری جامع در مقابل نوسانات متغیر محیطی را متجلی نمی‌کنند. علاوه بر این ابزارهای پیشرفته مانند تخصیص مقاصد بازپخش گذشته امکان استفاده‌ی کارآمد مبتنی بر اهداف نامشخص سیستمِ خودوفقی را دارا نیستند.

بر این اساس یک \textbf{شکاف علمی و عملی مشخص} برای ابداع چارچوبی مبرهن مطرح است که دارای ویژگی‌های زیر باشد:
۱) تسلط اثربخش در فضاهای وسیع و گسسته وفقی بدون استفاده از جداول سنتی مقیاس ناپذیر،
۲) فراگیری هوشمندانه برخط بدون اجبار استفاده از داده‌های آموزش دیده‌ی گذشته،
۳) مقابله سرسختانه با بهینه‌های محلی بدون از دست دادن همگرایی و افت در محیط تغییریافته (نیاز به شکل‌دهی استراتژی آموزش)، و
۴) قابلیت پیوستگی با نیازهای غیرخطی یک سیستم خودوفقی بدون تغییر پارادایم تحمیلی به مسائلی متمرکز به تعیین هدف.

پیشنهاد بنیادی ما در این رساله مبتنی بر ایجاد یک مکانیزم جدید تلفیقی تحت عنوان \textbf{یادگیری تقویتی عمیق با شکل‌دهی تصادفی پاداش (\lr{RS-DRL})} است تا خلأ یاد شده کاملاً پر گشته و گامی کاربردی جهت ارتقای استقلال عوامل کنترلی نرم‌افزارهای پیچیده میسر گردد.

\section{خلاصه فصل}

این فصل در آغاز، با ترسیم بنیان‌های نظری مفاهیم یک سیستم پایدار و مکانیزم \lr{MAPE-K} و تشریح ماهیت مدل‌سازیِ عدم‌قطعیت با ساختار زمان‌اجرا بنا گردید. به تبع آن ضمن تشریح کارکردهای یادگیری تقویتی متدهای عمیق پیوستگی استراتژی پیشنهادی، محدوده‌های فعلی در مقابله با چالش مقیاس‌پذیری نمایان شد. در خاتمه، بر پایه ارزیابیِ تحلیلی در سابقه کارهای تحقیقاتی پیشین مشخص گردید، فقدان یک مکانیزم هوشمند برای شکل‌دهی تصادفی برخط کاملا مشهود است که در فصل آتی با مدل سیستم و فرمول‌سازی مسئله پیش‌نیازهای ساخت آن تنظیم می‌گردد.

\newpage

% فصل ۳: مدل سیستم و فرمول‌بندی مسئله
\chapter{مدل سیستم و فرمول‌بندی مسئله}
\section{مقدمه}

برای توسعه یک راهکار مبتنی بر یادگیری تقویتی به منظور مدیریت عدم قطعیت در سیستم‌های خودوفقی، ابتدا باید یک مدل صوری از مسئله وفقی‌سازی (تطابق) بنا نهیم. این فصل، مدل سیستم و فرمول‌بندی ریاضیاتی که بستر رویکرد پیشنهادی ما (\lr{RS-DRL}) را تشکیل می‌دهد، ارائه می‌نماید. ما با معرفی \lr{DeltaIoT} -- یک نمونه بارز از شبکه‌های حسگر بی‌سیم که چالش‌های کلیدی فضاهای وفقی بزرگ و عدم قطعیت‌های محیطی را به صورت آشکار به نمایش می‌گذارد -- بحث را آغاز می‌کنیم.

سپس مسئله وفقی‌سازی را به عنوان یک فرآیند تصمیم‌گیری مارکوف (\lr{MDP}) فرمول‌بندی کرده و فضای حالت، فضای عمل و تابع پاداش را برای سناریوهای بهینه‌سازی تک‌هدفه و چندهدفه تعریف می‌کنیم. نهایتاً، مسئله یادگیری یک سیاست وفقی بهینه را که می‌بایست در شرایط متغیر زمان اجرا و به طور مداوم عمل کند، به صورت صوری بیان می‌نماییم.

\section{مطالعه موردی انگیزشی: سیستم \lr{DeltaIoT}}

\subsection{مرور کلی سیستم و منابع عدم قطعیت}

سیستم \lr{DeltaIoT} یک شبکه حسگر بی‌سیم (\lr{WSN}) با تعاملات پیچیده است که به صورت گسترده به عنوان یک محک استاندارد (\lr{Benchmark}) برای ارزیابی سیستم‌های خودوفقی به کار می‌رود \cite{RS_DRL_Article}. این شبکه بسته به نسخه مورد استفاده، مابین ۱۵ تا ۳۷ گره (سنسور) دارد. دو منبع اصلی عدم قطعیت در این سیستم در زمان اجرا بروز می‌یابند:
\begin{enumerate}
    \item \textbf{نویز و تداخل شبکه (\lr{Network Interference/Noise}):} تغییرات محیطی و مداخلات پیرامونی که موجب نوسانات پیش‌بینی‌نشده در کیفیت ارتباطات از \lr{-40dB} تا \lr{+15dB} می‌گردد.
    \item \textbf{تغییرپذیری بار پیام‌ها (\lr{Message Load Variability}):} نوسان در تعداد پیام‌های تولید شده توسط هر سنسور (بین ۰ تا ۱۰ پیام) با الگوهای انتقال متفاوت در طول زمان.
\end{enumerate}

همانطور که در مطالعات پیشین (و روند نشان داده شده در شکل ۲ مقاله مرجع) نمایان است، این عدم قطعیت‌ها باعث شکل‌گیری افت شدید عملکرد سیستم در طی چرخه‌های وفقی گشته و اهمیت حیاتی مکانیسم‌های تطابق پویا را برجسته می‌کنند.

\subsection{فضای وفقی و بازنمایی پیکربندی}

فضای وفقی (\lr{Adaptation Space}) دربرگیرنده مجموعه تمامی گزینه‌ها و پیکربندی‌های ممکنی است که سیستم می‌تواند در زمان اجرا اتخاذ نماید. در \lr{DeltaIoT}، سه نسخه متفاوت با مشخصاتی به شرح زیر طراحی گردیده است \cite{RS_DRL_Article}:
\begin{itemize}
    \item \textbf{\lr{DeltaIoTv1}:} شامل ۱۵ سنسور و دارای ۲۱۶ گزینه وفقی متمایز.
    \item \textbf{\lr{DeltaIoTv1.1}:} شامل ۱۶ سنسور و دربرگیرنده ۱۲۹۶ گزینه وفقی است.
    \item \textbf{\lr{DeltaIoTv2}:} نسخه وسیع‌تر این شبکه متشکل از ۳۷ سنسور که دارای فضای عظیم متشکل از ۴۰۹۶ تصمیم ممکن می‌باشد.
\end{itemize}

هریک از این گزینه‌های وفقی در حقیقت ترکیبی از پارامترهای مختلف هستند؛ نظیر سطوح توان انتقال انرژی (\lr{Transmission Power}) برای سنسورهای مشخص و توزیع نسبی مسیرهای شبکه‌ای (\lr{Link Distribution}). انتخاب یک ترکیب از این پارامترها، یک پیکربندی کاملاً جدید و واحد را شکل می‌دهد.

\section{مدل صوری برای تصمیم‌گیری خودوفقی تحت عدم قطعیت}

برای امکان‌پذیر کردن یادگیری تقویتی، نیاز به فرمول‌بندی مسأله تطابق شبکه در قالب یک فرآیند تصمیم‌گیری مارکوف (\lr{MDP}) داریم که نمایانگر چهارگانه $(S, A, R, P)$ است \cite{RS_DRL_Article}. 

\subsection{تعریف فضای حالت}

فضای حالت مدل، وضعیت فعلی و پویای سیستم را با اندازه‌گیری سه معیار حیاتی عملکرد بازتاب می‌دهد. به بیان رسمی:
$$S = \{(E, P, L) \mid E \in \mathbb{R}, P \in [0,100], L \in [0,100]\}$$
که در آن:
\begin{itemize}
    \item \textbf{انرژی مصرفی ($E$):} مقدار حقیقیِ نمایانگر انرژی کل مصرفی توسط تمام سنسورهای شبکه در چرخه پیشین است.
    \item \textbf{تلفات بسته‌های شبکه ($P$):} درصدی از پیام‌های تولید شده که در طول مسیر به مقصد (دروازه اصلی) از دست رفته‌اند (بین ۰ تا ۱۰۰ درصد).
    \item \textbf{تاخیر ارسال ($L$):} معیاری معادل زمان مورد نیاز جهت رسیدن اطلاعات از سنسورها به گیت‌وی (حداکثر بازه ۱۰۰).
\end{itemize}

این سه متغیر مقداری به این دلیل انتخاب شده‌اند که تصویر کاملی از نیازمندی‌های کیفی (\lr{QoS}) شبکه را در یک سیستم توزیع‌شده با منابع محدود ارائه کرده و ارزیابی صحیحی برای عامل تصمیم‌گیرنده (\lr{Agent}) فراهم می‌کنند.

\subsection{فضای عمل: گزینه‌های وفقی گسسته}

فضای عمل شامل تمامی استراتژی‌ها و پیکربندی‌های ممکن به صورت مجموعه‌ای از گزینه‌های گسسته تعریف می‌شود:
$$A = \{a_1, a_2, \ldots, a_n\}$$
همان‌طور که پیش‌تر اشاره شد، مقدار $n$ بسته به ابعاد شبکه \lr{DeltaIoT} برابر با ۲۱۶، ۱۲۹۶ یا ۴۰۹۶ در نظر گرفته می‌شود. هر عمل $a_i$ نمایانگر یک پیکربندی مشخص از پارامترهای $\{p_{i1}, p_{i2}, \ldots, p_{im}\}$ است که در آن هر پارامتر $p_{ij}$ می‌تواند از مجموعه مقادیر محدود $\{v_{ij1}, v_{ij2}, \ldots, v_{ijk}\}$ تبعیت کند. این پارامترها در واقع توان ارسال سنسورها، توزیع لینک‌ها و مسیر یابی هر سنسور را تغییر می‌دهند.

\subsection{فرمول‌بندی تابع پاداش}

محاسبه صحیح تابع پاداش برای هدایت سیستم به سمت مقاصد از پیش تعیین شده، عاملی کلیدی محسوب می‌شود. بر اساس ماهیت سیستم و اهداف، ممکن است با بهینه‌سازی‌های تک‌هدفه، چندهدفه و یا اهداف آستانه‌ای مواجه شویم.

\subsubsection{فرمول‌بندی پاداش چندهدفه}

برای متوازن‌سازی رفتار سیستم حین برخورد با اهداف متعارض در طول چرخه خودوفقی، یک تابع پاداش بهینه‌سازی چندهدفه ارائه می‌گردد:
$$R = w_E \cdot \left(\frac{E_{max} - E}{E_{max} - E_{min}}\right) + w_P \cdot \left(\frac{P_{max} - P}{P_{max} - P_{min}}\right) + w_L \cdot \left(\frac{L_{max} - L}{L_{max} - L_{min}}\right)$$
در این نگاشت، ضرایب وزن‌دهی ($w_E + w_P + w_L = 1$) میزان اهمیت هر هدف را نشان می‌دهند و هر ترم موجود با استفاده از کران‌های بالا و پایین هر معیار محدودیت‌ها را بین بازه $[0, 1]$ نرمال‌سازی می‌کند.

\subsubsection{فرمول‌بندی پاداش تک‌هدفه}

زمانی که اولویت طراحی تنها برای غلبه بر یک مشکل واحد تعریف گردیده (مانند بهینه‌سازی صرفاً مصرف انرژی)، تابع فوق بدین شکل ساده‌سازی می‌گردد:
$$R = \left(\frac{E_{max} - E}{E_{max} - E_{min}}\right) \cdot w_E$$
در این معادله بقیه وزن‌ها معادل صفر در نظر گرفته شده‌اند که سبب هدایت بدون واسطه سیستم برای ارتقای سریع ویژگی هدف‌گذاری شده می‌شود.

\subsubsection{پارامترهای پاداش برای اهداف آستانه‌ای (\lr{TTS} و \lr{TT})}

برای مسائلی که نیازمند کنترل مطلق اهدافی در سطح استثناهای کیفیِ قابل قبول (\lr{Thresholds}) هستند (که در جدول ۲ مقاله با نام‌های \lr{TTS} و \lr{TT} نام برده شده‌اند)، تخصیص مقادیر پاداشِ مشروطی به صورت دودویی (\lr{Binary}) اعمال می‌گردد:
$$R = \begin{cases} 
1 & \text{اگر } P < \theta_P \text{ و } L < \theta_L \\
0 & \text{در غیر این صورت}
\end{cases}$$
که در آن مقادیر مرزی $\theta_P$ و $\theta_L$ به عنوان آستانه‌های از پیش تعیین شده معرفی شده‌اند. جدول \ref{tab:deltaiot_goals} خلاصه مقادیر \lr{TTS} و \lr{TT} (تنظیمات اهداف آستانه‌ای و ترکیبی) را پیرامون معماری‌ این نسخه‌ها نمایش می‌دهد:

\begin{table}[h!]
\centering
\caption{اهداف کیفی و آستانه‌ای (\lr{TTS} و \lr{TT}) سیستم \lr{DeltaIoT} \cite{RS_DRL_Article}}
\label{tab:deltaiot_goals}
\begin{tabular}{|c|c|c|c|c|}
\hline
\textbf{نسخه \lr{DeltaIoT}} & \textbf{تنظیمات هدف} & \textbf{تلفات بسته (\lr{Packet Loss})} & \textbf{تاخیر (\lr{Latency})} & \textbf{انرژی (\lr{Energy})} \\ \hline
\lr{DeltaIoTv1} & \lr{TTS} & کمتر از ۱۰٪ ($<10\%$) & کمتر از ۵٪ ($<5\%$) & کمینه‌سازی \\ \hline
\lr{DeltaIoTv1} & \lr{TT} & کمتر از ۱۰٪ ($<10\%$) & - & کمینه‌سازی \\ \hline
\lr{DeltaIoTv2} & \lr{TTS} & کمتر از ۱۰٪ ($<10\%$) & کمتر از ۵٪ ($<5\%$) & کمینه‌سازی \\ \hline
\end{tabular}
\end{table}

\section{بیان مسئله: یادگیری سیاست وفقی بهینه با عملیات مداوم در زمان اجرا}

با در نظر داشتن فرمول‌بندی \lr{MDP} در نظر گرفته شده، مسأله نهایی ما یادگیری روال سیاست وفقی بهینه $\pi^*: S \rightarrow A$ است. انتظار سیستم از این سیاست به حداکثر رسانی ارزش پاداش‌های تجمعی تنزیل‌یافته (\lr{Cumulative Discounted Reward}) است:
$$\pi^* = \arg\max_{\pi} \mathbb{E}\left[\sum_{t=0}^{\infty} \gamma^t R(s_t, a_t) \mid \pi\right]$$
در معادله فوق، مقدار $\gamma \in [0,1)$ ضریب تنزیل را معرفی می‌نماید. عملکرد موثر این سیاست بهینه باید همواره به طور مداوم و لاینقطع در زمان اجرا عمل کند و مکرراً با چالش‌های زیر روبرو گردد:
\begin{enumerate}
    \item \textbf{عدم قطعیت محیطی:} بار متغیر پیام‌ها و همچنین تداخل‌های ناپایدار و اختلال متناوب خطوط.
    \item \textbf{فضای عمل گسسته گسترده:} مدیریت فضای مقیاس‌پذیری که بسته به نسخه شبکه می‌تواند تا سقف ۴۰۹۶ تصمیم ممکن وسیع باشد.
    \item \textbf{پویایی غیرایستا (\lr{Non-stationary dynamics}):} نیاز به تطبیق مداوم و عدم کاهش سرعت واکنش به علت تغییرات شرایط از یک چرخه وفقی به چرخه بعدی (تخریب و کاهش عملکرد نشان داده شده در شکل ۲ مقاله مرجع).
\end{enumerate}

\section{خلاصه فصل}

فرآیند صورت‌بندی و مدل‌سازی در این فصل، زمینه‌ساز توسعه راهکار عملی این رساله گشته است. ما در این بخش ابتدا مطالعه موردی سیستم \lr{DeltaIoT} را شرح داده و ابعاد گوناگون فضای وفقی گسسته‌ی آن را تشریح نمودیم. در ادامه این مسأله با تمرکز به مدل تصمیم‌گیری مارکوف ساختاریافته شامل فضای وضعیت شبکه، اعمال تصمیم‌سازی و توابع پاداش‌های چندگانه برای بهینه‌سازی فرموله گردید.

تکمیل این چارچوب رسمی، قابلیت و ظرفیت یادگیری استراتژی‌های وفقی و منطبق با متد یادگیری تقویتی را فراهم می‌آورد. این پایه‌گذاری با اهمیت به ویژه به واسطه توصیف اهداف مشخص، در فصل آتی (به عنوان طراحی متد ابتکاری \lr{RS-DRL}) زمینه‌ساز توسعه راهکارِ خروج از بهینه‌های محلی با بهره‌گیری از شکل‌دهی تصادفیِ پاداش برای مدیریت عدم قطعیت خواهد بود.

\newpage

% فصل ۴: روش پیشنهادی \lr{RS-DRL}
\chapter{روش پیشنهادی \lr{RS-DRL}: معماری و الگوریتم}
\section{مقدمه}

در فصل پیشین نشان دادیم که مدل‌های مرسوم یادگیری تقویتی، توانایی لازم برای پاسخ‌گویی به فضاهای بسیار بزرگ و گسسته تطابق را در زمان اجرای سیستم‌های خودوفقی ندارند و رویکردهای موجود اغلب در قله‌های محلی (\lr{Local Optima}) متوقف می‌شوند. همچنین آموزش مجدد مدل در زمان اجرا فرآیندی پرهزینه است و اتلاف وقت قابل توجهی برای یافتن گزینه‌های بهینه می‌طلبد. برای غلبه بر این محدودیت‌ها، سیستم نیازمند رویکردی است که بتواند با کمترین میزان نیاز به داده‌های از پیش‌تعیین‌شده و برچسب‌دار، تجربیات پیشین خود را ارزیابی و استراتژی‌های جدیدی را در زمان اجرا به کار گیرد.

این فصل رویکرد پیشنهادی رساله یعنی مدیریت عدم قطعیت مبتنی بر مدل ترکیبی «یادگیری تقویتی عمیق با شکل‌دهی تصادفی پاداش» یا به اختصار \lr{RS-DRL} را معرفی می‌نماید. ما در این فصل به صورت جامع، معماری مبتنی بر حلقه کنترل \lr{MAPE-K} که با الگوریتم‌های هوش مصنوعی تقویت شده است را توجیه خواهیم کرد و نشان می‌دهیم چگونه یک عامل با مدیریت بهینه و شکل‌دهی مجدد تجربیات ناموفق خود، می‌تواند سرعت همگرایی را افزایش داده و راه‌حل‌های بهینه‌سراسری را حتی در فضاهای بزرگ تطبیقی کشف نماید.

در ادامه، ابتدا مروری کلی بر ایده و روش پیشنهادی ارائه خواهد شد. سپس معماری یکپارچه سیستم و فرآیند ۹ مرحله‌ای در چرخه \lr{MAPE-K} شرح داده می‌شود. در نهایت، تمرکز اصلی بر روی جزئیات الگوریتم \lr{RS-DRL} و مکانیزم بدیع تصادفی‌سازی پاداش‌ها خواهد بود و مقایسه‌ای جامع بین این روش و بازپخش تجربه با نگاه به گذشته (\lr{HER}) انجام می‌گیرد.

\section{مرور کلی بر روش پیشنهادی}

روش \lr{RS-DRL} یک متدولوژی نوین است که بر پایه ارتقای الگوریتم سنتی شبکه‌های دیپ‌کیو (\lr{DQN}) طراحی شده تا فرآیند بازآموزی سیستم را، بدون نیاز به دانش قبلی یا تکیه بر آموزش‌های آفلاین دشوار، بهبود بخشد. ایده اصلی این روش، تعبیه یک مکانیزم شکل‌دهی پاداش بدیع (\lr{Reward Shaping}) در طول فاز حافظه بازپخش تجربه (\lr{Experience Replay}) می‌باشد.

نوآوری کلیدی نهفته در روش پیشنهادی آن است که به جای استفاده‌های یکنواخت و مرسوم از تمامی تجربیات گذشته، سیستم، زیرمجموعه‌ای از تجربیات شکست‌خورده (که در آنها مقدار پاداش $r_i=0$ بوده است) را استخراج کرده و به صورت تصادفی آنها را به عنوان تجربیات موفق (با تخصیص پاداش خوش‌بینانه $r_i=1$) بازطراحی و شکل‌دهی مجدد می‌نماید. این ترفند استراتژیک، عامل یادگیرنده را تشویق می‌کند تا به جای گیر افتادن در سیاست‌های محلی ضعیف گذشته، با جسارتِ بیشتر کاوش (\lr{Exploration}) بیشتری انجام دهد و در نهایت با افزایش قدرت تعمیم‌پذیری، استراتژی‌های جدیدتری را در مواجهه با شرایط عدم قطعیت پیدا کند.

\section{معماری سیستم: یکپارچه‌سازی \lr{RS-DRL} با حلقه \lr{MAPE-K}}

\subsection{سیستم مدیریت‌شده و سیستم مدیریت‌کننده}

همان‌طور که در فصل‌های پیشین به تفصیل بیان شد، جداسازیِ منطقِ کسب‌وکار از منطق انطباق یکی از اصول پایه‌ای در سیستم‌های خودوفقی به شمار می‌رود. بر این اساس، معماری پیشنهادی ما از دو بخش مجزا تشکیل گردیده است:
\begin{enumerate}
    \item \textbf{سیستم مدیریت‌شونده (\lr{Managed System}):} این بخش همان سیستم نرم‌افزاری پایه‌ای است که نیازمند فرآیند وفقی‌سازی (تطبیق) می‌باشد. در روند این رساله، پلتفرم \lr{DeltaIoT} به عنوان یک شبکه حسگر بی‌سیم، تحت عنوان سیستم مدیریت‌شونده در نظر گرفته شده است که داده‌های محیطی زمانِ اجرا را تولید می‌نماید.
    \item \textbf{سیستم مدیریت‌کننده (\lr{Managing System}):} این بخش شامل مکانیزم حلقه کنترل و بازخورد \lr{MAPE-K} است که به صورت عمیقی با عامل یادگیری تقویتی \lr{RS-DRL} ادغام یافته تا مدیریتِ پویای تصمیماتاتخاذ شده و حل مسئله بهینه‌سازی را بر عهده گیرد. 
\end{enumerate}

\subsection{تعامل دقیق حلقه وفقی ۹ مرحله‌ای}

ادغام یک مدل یادگیری تقویتی مانند \lr{RS-DRL} درون یک سیستم حیاتی، نیازمند یک جریان یکپارچه و امن از دریافت داده، ارزیابی وضعیت، به‌روزرسانی مدل، تأیید صوری و نهایتاً اجرای گزینه‌ها است. این فرآیند از طریق ۹ مرحله تعاملی بین مؤلفه‌های مختلف در حلقه مدیریت مشخص می‌گردد:

\begin{enumerate}
    \item \textbf{جمع‌آوری داده‌ها توسط مؤلفه پایش (\lr{Monitor}):} در گام نخست، این مؤلفه داده‌های زمان اجرای سیستم نظیر خوانش‌های سنسوری و لاگ‌های سیستم را به شکل مداوم جمع‌آوری می‌نماید.
    \item \textbf{ارسال داده‌ها جهت آموزش مدل:} اطلاعات جمع‌آوری شده شامل وضعیت رخ‌داده، پیامدهای استراتژی‌های پیشین و توابع پاداش محاسبه شده، برای بازآموزی عامل هوشمند، مستقیماً به ماژول \lr{RS-DRL} تزریق می‌گردد.
    \item \textbf{ذخیره‌سازی مدل آموزش‌دیده در پایگاه دانش (\lr{Knowledge}):} پس از اتمام فاز بازآموزی و اجرای مکانیزم بازپخش تجربه با ویژگی‌های شکل‌دهی تصادفی پاداش‌ها، مقادیر وزن‌ها و سیاست‌های به‌روز شبکه عصبی بر روی مخزن دانش مرکزی سیستم استقرار می‌یابد.
    \item \textbf{انتقال همزمان داده‌ها به تحلیل‌گر (\lr{Analyzer}):} مؤلفه پایش در موازات عملیات یادگیری، داده‌های شرایط کنونی را جهت بررسی انطباق‌ها به سمت مؤلفه تحلیل‌گر نیز هدایت می‌کند.
    \item \textbf{آشکارسازی نقض اهداف کیفی:} تحلیل‌گر، انطباقِ وضعیت جاری سیستم با ضمانت‌ها و اهداف از پیش تعیین‌شده را ارزیابی می‌کند. در صورت مشاهده نقض اهداف کیفی (نظیر افت طولانی منابع یا اختلال)، بلافاصله هشداری را برای مؤلفه برنامه‌ریز (\lr{Planner}) ارسال می‌دارد.
    \item \textbf{استخراج مدل توسط انتخاب‌گر گزینه وفقی:} با دریافت هشدار از سوی تحلیل‌گر، ماژول پیش‌بینی کننده گزینه وفقی (\lr{Adaptation Option Predictor}) که در قلب مؤلفه برنامه‌ریز قرار دارد، آخرین نسخه آموزش‌دیده و به‌روز از مدل \lr{RS-DRL} را از بخش دانش برای استفاده استخراج می‌کند.
    \item \textbf{پیش‌بینی و انتخاب گزینه‌ی بهینه:} با فراهم شدن مدل عملیاتی، انتخاب‌گر وضعیت جدید شبکه را به شبکه عصبی وارد نموده و استراتژی و گزینه تطبیقی پیش‌بینی شده در توابع کیفیت را انتخاب می‌کند.
    \item \textbf{دریافت دستورالعمل‌های گزینش از پایگاه دانش:} مؤلفه مجری (\lr{Executor}) برنامه تاییدشده را از برنامه‌ریز اخذ کرده و برای اطلاع از مختصاتِ عملیاتی، تنظیمات سطح پایین و جزئیات دقیق اجرای آن، مجدداً به مخزن دانش اطلاعات رجوع می‌نماید.
    \item \textbf{اجرای قطعی فرامین وفقی:} در گام پایانی، مؤلفه مجری، فرمان تطبیقی اخذ شده (مانند تغییر توان ارسال و تخصیص مسیرهای جدید برای گره‌ها در \lr{DeltaIoT}) را بر روی ساختار سیستمِ مدیریت‌شده جاری می‌سازد.
\end{enumerate}

\section{الگوریتم \lr{RS-DRL} و مکانیزم تصادفی‌سازی}

این بخش، هسته الگوریتمی که سبب برتری کامل سیستم در حفظ کیفیت در فضاهای متغیر و بزرگ شده را تبیین می‌نماید.

\subsection{الگوریتم اصلی: \lr{DQN} برای مدیریت حلقه‌های وفقی جریانی}

برای شبکه‌های عظیم و گسسته‌ای که دارای پویایی بالا می‌باشند، ساختار الگوریتمی شبکه‌های \lr{Q} عمیق (\lr{DQN}) ویژگی‌های مطلوبی را فراهم می‌سازد. انتخاب این چارچوب یادگیری تقویتی بر اساس چهار برهان مستحکم صورت گرفته است: 
۱) کارایی بنیادین: تسلط و پردازش سریع فضاهای عمل گسسته به شیوه‌ای کارآمد بدون از دست دادن زمان.
۲) آموزندگی بهره‌ور: دارا بودن حافظه بازپخش که سبب کارایی در نمونه‌برداری از تجربیات رویدادهای تاریخی (\lr{Sample Efficiency}) می‌گردد.
۳) ایستادگی و ثبات: بهره‌گیری از شبکه‌های تثبیت‌شده هدف (\lr{Target Networks}) که ثبات اهداف و کمینه‌کردن نوسانات شیب‌یابی در سیستم‌های متغیر را تضمین می‌نمایند.
۴) یادگیری از تجارب تاریخی: فرآیند یادگیری مبتنی بر مدل‌های غیرسیاستی (\lr{Off-Policy}) که قابلیت تغذیه آفلاین و استخراج الگو از سوابق تجربیات متوالی را برای عامل مهیا می‌سازد. 

\subsection{ایده اصلی: شکل‌دهی تصادفی پاداش در طول بازپخش تجربه}

متمایزترین عامل نوآوری سیستم در این است که برخلاف مکانیسم حافظه‌های سنتی که منحصراً نمونه‌های گذشته را به صورت نمونه‌برداری یکنواخت (\lr{Uniform Sampling}) در فاز آموزش لحاظ می‌نمایند، مدلِ معرفی شده، زیرمجموعه‌ای تصادفی از تجربیاتِ تاریخچه‌دار که مستقیماً در انجام هدفِ کیفی شکست خورده‌اند (با اعطای مقدار $r=0$) را استخراج نموده، تغییر داده (برچسب مجدد) و وانمود می‌کند این اعمال در وضعیت‌های پیشین با موفقیت‌آمیز همراه بوده‌اند (تخصیص $r=1$). 

این فرآیند به صورت مطلق و کور صورت نمی‌پذیرد؛ بلکه توسط عاملی با عنوان \textbf{فاکتور شکل‌دهی $\rho$} (\texorpdfstring{$\rho$}{rho}) کنترل شده تا میزان و نرخ تغییرات تصادفی را تنظیم نماید. فلسفه و بصیرت طراحی (\lr{Intuition}) مستتر در این تکنیک این است که تخصیص امتیازات مثبت و ثابت به استثناهای ناموفقِ پیشین، عاملی محیط‌شناسانه ایجاد می‌نماید تا به جای اجتناب از بخش‌های مبهم حالت‌ها به دلیل شکست قبلی، مجدداً آن‌ها را در دوره‌های آتی با انگیزه اکتشافه مواجه ارزیابی کرده و بدین ترتیب با غلبه بر ترس از افت پاداش، راه‌حل‌های جهانی و ناشناخته‌تر را کشف کند. این امر از فروافتادن مدل به تله‌های بهینه محلی که ناشی از عدم بررسی مجدد انتخاب‌های پرخطر گذشته است، جلوگیری می‌نماید.

\subsection{تابع شکل‌دهی پاداش (الگوریتم ۲)}

مبانی دقیق کارکرد نظری سیستم در خصوص اصلاح پاداش تجربه‌ها، در الگوریتم \ref{alg:reward_shaping} شرح داده شده است. این روند پردازشی با دریافت مجموعه‌ای از نمونه‌برداری‌های تصادفی در قالب یک مینی‌بچ (\lr{Minibatch}) و اعمال فاکتور کنترل‌کننده $\rho$ آغاز می‌گردد.

در این رویه پردازشی، مکانیزم در گام‌های ۱ تا ۳، ابتدا مجموعه عملگرهای مردود شده از منظر پاداش را بررسی می‌کند. اگر تعداد رویدادهای مردود شده حدنصاب لازم محاسبه شده از طریق ضریب $\rho$ را داشته باشند، زیرمجموعه‌ای تصادفی شامل $K$ مقدار از این دسته، گزینش شده و مقدار متغیر پاداش آن‌ها به عدد ۱ تغییر می‌یابد (\lr{Reshape}) و نهایتاً دسته داده‌های به‌روزشده بازگردانده خواهند شد. استفاده از یک مقدار ثابت خوش‌بینانه نظیر استراتژی (تخصیص $1$) یک مسیر مستقیم و غیر‌وابسته را مهیا نموده که با تغییر رفتار فرآیند یادگیری انسان‌ها شباهت دارد.

\begin{algorithm}[H]
\caption{شکل‌دهی پاداش (\lr{Reward Shaping})}
\label{alg:reward_shaping}
\begin{algorithmic}[1]
\renewcommand{\algorithmicrequire}{\textbf{ورودی:}}
\renewcommand{\algorithmicensure}{\textbf{خروجی:}}
\REQUIRE مینی‌بچ $B$، فاکتور ضریبی شکل‌دهی $\rho$
\STATE $N \leftarrow |B|$
\STATE $K \leftarrow \rho \times N$ \COMMENT{محاسبه تعداد تجربیاتی که باید شکل‌دهی شوند}
\STATE شناسایی تجربیات ناموفقِ پیشین به شرط داشتن مقدار $r_i = 0$
\IF{تعداد کل تجربیات ناموفق مینی‌بچ $\geq K$}
    \STATE انتخاب تصادفی متعادل از میان $K$ تجربه ناموفق
    \FOR{هر تراکنش انتخاب شده نظیر $(s_i, a_i, 0, s_{i+1})$}
        \STATE $r_i \leftarrow 1$ \COMMENT{شکل‌دهی یک شکست به موفقیت تصنعی}
    \ENDFOR
\ELSE
    \STATE تمام تجربیات ناموفق حاضر را با پاداش ثابت ۱ شکل‌دهی کن
\ENDIF
\RETURN جریان مینی‌بچ به‌روزشده $B$
\end{algorithmic}
\end{algorithm}

\subsection{الگوریتم جامع و کلان \lr{RS-DRL} (الگوریتم ۱)}

الگوریتم \ref{alg:rs_drl} روند کامل سیستمِ هوشمند از ساختار شبکه عصبی تا حلقه‌های تعاملی با مکانیزمِ معرفی‌شدۀ فوق را تبیین می‌نماید. 

خطوط ۱ تا ۲ فرآیند راه‌اندازی، تثبیت پارامترهای پایه‌ای و تخصیص شیب‌های اولیه برای شبکه اصلی و شبکه هدف را مستند می‌نمایند. خطوط ۳ تا ۵ نشان‌دهنده حلقه‌های اجرای تو در توی اپیزودهای پیوسته زمان اجراست که با انتخاب‌های اتخاذی از سوی سیاست حریصانه خط ۶ (\lr{$\epsilon$-greedy}) هنجارهای تطبیقی را صورت‌بندی می‌نمایند. کلیه این تجربه‌های محیطی اجرا شده نیز در حافظه‌ی سیستم افزوده می‌گردند (خطوط ۷-۸). خط ۹ نمونه‌گیری اتفاقی (\lr{Sampling}) را اعمال داشته، و پس از آن الگوریتم شماره ۲ (تغییرشکل‌پاداش) بر مدل فراخوانده می‌شود (خط ۱۱) تا ساختارها را تنظیم مجدد کند. نهایتاً با یافتن اهداف درختیِ خطاها، وزن‌دهی گرادیان جهت ارتقای دقت مدل، در خطوط ۱۲ تا ۱۴ بروز می‌گردد.

\begin{algorithm}[H]
\caption{الگوریتم یکپارچه یادگیری \lr{RS-DRL}}
\label{alg:rs_drl}
\begin{algorithmic}[1]
\STATE راه‌اندازی شبکه Q اصلی و تعیین وزن‌های تصادفی $\theta$
\STATE راه‌اندازی و تخصیص شبکه تثبیت شونده هدف با وزن‌های $\theta' = \theta$
\FOR{هر اپیزود تعاملی}
    \STATE دریافت حالت اولیه محیط $s$
    \WHILE{حلقه اپیزود تمام نشده است}
        \STATE انتخاب اقدام انطباقی $a$ با استفاده از سیاست $\epsilon$-گریدی
        \STATE اجرای اقدام $a$ بر روی محیط، و دریافت پاداش $r$ و حالت غایی حاصله $s'$
        \STATE ذخیره‌سازی این تراکنش به فرم $(s, a, r, s')$ در حافظه بازپخشِ $\mathcal{D}$
        \STATE اخذ نمونه‌گیریِ یک مینی‌بچ رندوم $B$ از $\mathcal{D}$
        \IF{شرایط لازم برای بروزرسانیِ شکل‌دهی مهیا باشد}
            \STATE $B \leftarrow \text{RewardShaping}(B, \rho)$ \COMMENT{فراخوانی الگوریتم \ref{alg:reward_shaping}}
        \ENDIF
        \STATE برآورد هدف تقربی $Q$ برای نمونه‌های داخل مینی‌بچ با شبکه ثابت
        \STATE اجرای کاهنده نزول گرادیان روی شبکه اصلی برای کمینه‌سازی اختلافات
        \STATE انتقال و به‌روزرسانی وزن‌های شبکه هدف $\theta'$ به صورت درصدی
        \STATE تخصیصِ وضعیتِ فعلی: $s \leftarrow s'$
    \ENDWHILE
\ENDFOR
\end{algorithmic}
\end{algorithm}

\subsection{مقایسه با بازپخش تجربه با نگاه به گذشته (\lr{HER})}

استراتژی‌های متداولی پیرامون کار با تجربیاتِ مفقودِ سیستم در حوزه هوش مصنوعی توسعه یافته که مشهورترین مدل‌ها متأثر از استراتژی بازپخش تجربه با نگاه به گذشته یا \lr{HER} (و نمونه اخیر \lr{Soft HER}) می‌باشند. به منظور شفاف‌سازی ویژگی‌های کلیدیِ این طراحی و اثبات برتریِ آن در محیط‌هایِ خودسازمانده، جدول \ref{tab:comparison_her} جنبه‌های کلیدی تمایز این رویکردها را گردآوری نموده است.

\begin{table}[h!]
\centering
\caption{تحلیل مقایسه‌ای ویژگی‌های \lr{RS-DRL} با راهکارهای مبتنی بر \lr{HER}}
\label{tab:comparison_her}
\begin{tabular}{|c|c|c|c|c|}
\hline
\textbf{متد یا الگو} & \textbf{نیاز به تعریف صریح هدف} & \textbf{وابسته به دانش دامنه‌ای} & \textbf{خاستگاهِ استراتژی شکل‌دهی} & \textbf{دامنه کاربرد و استقرار} \\ \hline
\lr{HER} & بله (الزامی) & دارای تسلط قطعی به محیط & برچسب‌زنی مجدد و صریح اهداف (\lr{Goal relabeling}) & کاربرد محدود به محیط‌های غایت‌گرا \\ \hline
\lr{Soft HER} & بله (الزامی) & تسلط نسبی به پارامترهای مدل & برچسب‌زنی ملایم و تصادفی برای مقاصد & مقاصد کاملاً مشخص و معین در پیوسته \\ \hline
\lr{RS-DRL} & \textbf{خیر} & \textbf{اساساً نامتکی و بدون نیاز} & \textbf{تخصیص تصادفی مقدار موفقیت‌آمیز} & \textbf{محیط‌های مبهم سیستم‌های خودوفقی} \\ \hline
\end{tabular}
\end{table}

همانطور که ارزیابی گردید، الگوریتم‌های برگرفته از چارچوب \lr{HER} با وجود توانایی‌های شگرفِ جستجویی خود، نیازمندِ وابستگی‌های بسیار مهمی می‌باشند. در این متدها محقق باید برای الگوریتم، مسیرِ اهدافِ صریح و مقاصد کاملاً تعریف‌شده‌ای را برنامه‌ریزی نماید تا فرآیند برچسب‌گذاریِ مجدد (\lr{Goal Relabeling}) به درستی محقق شود. در مقابل استراتژی \lr{RS-DRL} اساساً یک استراتژیِ خودکفا، مقاوم و مستقل است که بدون دخالت انسانی و در نبود دانش دامنه‌ای، نیازهای خود را تنها متکی به پارامترهای استخراج شده از رویکردهای احتمالی پیدا می‌کند و بنابراین مناسب‌ترین مدل برای شبکه‌های نامتقارن، محیط‌های نامشخص و پویایی‌های غیرایستای سیستم‌های خودوفقی قلمداد می‌گردد.

\section{جزئیات کلیدی پیاده‌سازی برای بازتولیدپذیری}

\subsection{ساخت حالت از داده‌های پایش شده}
در سیستمِ مدل‌سازی‌شده برای اینکه متغیرهای پراکنده به الگویی معنادار جهت ارزیابی در زمانِ اجرا بدل شوند باید وضعیت سیستم را تحت قالب فضای حالت مهیا ساخت. مطابق با قواعد استخراج شده، داده‌های خام دریافتی مستقیماً به الگویی به فرمت زیر تبدیل می‌شوند: 
$$S = \{(E, P, L) \mid E \in \mathbb{R}, P \in [0,100], L \in [0,100]\}$$
ایجاد ساختار وضعیت وابسته به سه ورودیِ اساسی بوده است:
\begin{enumerate}
    \item \textbf{معیار انرژی ($E$):} استنتاج شده از مقدار انرژی مصرف شده (به صورت مقادیر حقیقی شناور).
    \item \textbf{معیار پکت‌لاس ($P$):} ضریبِ درصدیِ بسته‌های مفقود و نابود شده، مقیاس‌بندی شده از بازه ۰ تا ۱۰۰.
    \item \textbf{معیار تاخیر انتقال ($L$):} درصد بالاترینِ سیگنال‌های دریافتی از مقیاس شبکه، بیانگر مدت زمان تأخیر از مبدأ به گره اصلی (تا بازه محدود ۱۰۰).
\end{enumerate}

\subsection{اعتبارسنجی زمان اجرا با \lr{UPPAAL-SMC} و \lr{ActivFORMS}}
اعمال استراتژی‌های جدید تطبیقی، به دلیل تبعیت از متدهای یادگیری که گاهاً با عدم قطعیت‌هایی روبرو هستند، به تضمین ایمنی و تایید صلاحیت صوری (\lr{Formal Verification}) نیاز دارد. به این دلیل، شبکه‌های پیش‌بینی شدهِ از سوی عامل یادگیرنده، بلافاصله در کنترلر \lr{UPPAAL-SMC} با بهره‌گرفتن از شبکه‌های مدل‌سازیِ اتوماتای زمان‌دار، از طریق اعتبارسنجی کنترل آماری مدل (\lr{Statistical Model Checking}) تحلیل و در صورت رعایت ایمنی‌های الزامی، توسط ماشین مجازی واسط \lr{ActivFORMS} جهت اجرا روی شبکه واقعی صادر می‌گردند.

\subsection{انتخاب و بهینه‌سازی ابرپارامترها}
جهت استقرار بالاترین راندمان سیستم شبکه‌عصبی، فرآیند سیستماتیک جستجوی شبکه‌ای (\lr{Grid Search}) برای کاوش جامع درون فضای تنظیمات تدارک دیده شده است. تنظیماتی نظیر تعیین نرخ یادگیری بهینه ($\eta$)، عامل تنزیل پاداش‌های بلند‌مدت ($\gamma$)، نرخ اکتشاف در فاز کاوشی ($\epsilon$) و البته سرنوشت‌سازترین شاخصِ رساله یعنی فاکتور شکل‌دهی مقدار تجربی تصادفی ($\rho$)، همگی بر پایه مقایسه سرعتِ همگرایی کیفیت خطِ راهکارها و تضمین ثباتِ کیفیت سیستم کالیبره و در مدل آزمایشگاهی تثبیت گردیده‌اند.

\section{خلاصه فصل}

در این فصل رساله با یک گذر ساختاریافته از نظریات تا گام‌های عملیاتی، معماری تمام عیاری جهت استقرار هوش مصنوعی تصمیم‌ساز و مدل‌سازی در بطن یک شبکه مدیریت خودوفقی، تبیین و مستدل گردید. یکپارچگیِ حلقه بازخورد معروفِ \lr{MAPE-K} با یادگیریِ تقویتی این ظرفیت را به سیستم مدیریتی اعطا می‌کند تا گره‌های شبکه را بدون کنترل انسانی مهار و تنظیم نماید.
در این راستا با معرفیِ شفافِ قالب‌های الگوریتمی وابسته، نوآوریِ مکانیزمِ اختصاصی این رساله در خصوص تخصیص ارزش مثبت به تجربیات ناموفق سیستمِ ثبت‌شده در حافظه‌ی بازپخش (\lr{RS-DRL}) و نقش بی‌نظیر فاکتور کنترلیِ شکل‌دهی تشریح شد و نشان داده شد که استراتژیِ ابتکاری نه تنها نیازمند دانش از قبلِ دامنه‌ی نرم‌افزار نیست، بلکه در مواجهه با چالش‌های سیستم‌های نوین نظیر \lr{HER} دارای برتری قطعیِ اجرایی جهت ممانعت از همگرایی به استثناهای محلی (\lr{Local Optima}) می‌باشد. در فصل آتی به صورت عملی شواهد این برتریِ محاسباتی و زمان اجرا در مقایسه دقیق متدهای دیگر ارائه خواهد شد.

\newpage

% فصل ۵: ارزیابی تجربی
\chapter{ارزیابی تجربی}
\section{مقدمه}

\section{اهداف ارزیابی و سوالات تحقیق}

\section{تنظیمات آزمایشگاهی}
\subsection{نمونه‌های مطالعه موردی: \lr{DeltaIoTv1}، \lr{DeltaIoTv1.1} و \lr{DeltaIoTv2}}
\subsection{مجموعه داده‌ها: ردپاهای محیطی و پارامترهای رانش}
\subsection{پلتفرم پیاده‌سازی: \lr{Gymnasium} و \lr{Stable-Baselines3}}

\section{روش‌های پایه و مقایسه}
\subsection{روش \lr{Epsilon-Greedy DQN}}
\subsection{روش \lr{RS-DRL} به همراه \lr{HER / Soft HER}}
\subsection{روش‌های \lr{DLASER}، \lr{DLASER+} و \lr{ML4EAS}}
\subsection{روش مرجع جستجوی کامل (\lr{Exhaustive Search})}
\subsection{توجیه انتخاب روش‌های پایه: محدودیت‌ها و پروتکل‌های مقایسه}

\section{معیارهای ارزیابی}
\subsection{عملکرد یادگیری: عملکرد مجانبی، زمان تا آستانه و عملکرد کل}
\subsection{عملکرد سیستم: از دست رفتن بسته، تاخیر، مصرف انرژی}

\section{نتایج و تحلیل‌ها}
\subsection{منحنی‌های یادگیری و رفتار همگرایی}
\subsection{مقایسه عملکرد به ازای هر هدف}
\subsection{ارزیابی مقایسه‌ای در برابر تمام روش‌های پایه}
\subsection{مطالعه ابلیشن (\lr{Ablation}): حساسیت به فاکتور شکل‌دهی (\texorpdfstring{$\rho$}{rho})}
\subsection{تحلیل مقیاس‌پذیری (از ۲۱۶ به ۴۰۹۶ گزینه وفقی)}

\section{اعتبارسنجی آماری (آزمون‌های \lr{t} و مقادیر \lr{p})}

\section{خلاصه فصل}

\newpage

% فصل ۶: بحث، کارهای آینده و نتیجه‌گیری
\chapter{بحث، کارهای آینده و نتیجه‌گیری}
\section{مقدمه}

\section{بحث در مورد یافته‌های کلیدی}
\subsection{نظریه‌پردازی مکانیزم: چرا \lr{RS-DRL} کار می‌کند؟}
\subsection{تحلیل تعاملات و مصالحه‌ها (به عنوان مثال، تاخیر در مقابل از دست رفتن بسته)}
\subsection{پیامدهای مهندسی سیستم‌های خودوفقی}

\section{محدودیت‌ها و تهدیدات اعتبار}
\subsection{اعتبار داخلی، خارجی و سازه}
\subsection{محدودیت‌های رویکرد \lr{RS-DRL} (وابستگی به \texorpdfstring{$\rho$}{rho}، فضاهای عمل گسسته)}

\section{مسیرهایی برای تحقیقات آینده}
\subsection{به سوی استراتژی‌های شکل‌دهی وفقی (تنظیم خودکار \texorpdfstring{$\rho$}{rho})}
\subsection{گسترش \lr{RS-DRL} به دامنه‌های پیوسته و حساس به ایمنی}
\subsection{اعتبارسنجی بین‌دامنه و تحلیل‌های نظری}

\section{نتیجه‌گیری}
\subsection{خلاصه بیان مسئله و راهکار پیشنهادی}
\subsection{مرور مشارکت‌ها و شواهد تجربی}
\subsection{سخنان پایانی}

\newpage

% کتابنامه و پیوست‌ها
\pagestyle{empty}
{
\onehalfspacing
\bibliographystyle{ieeetr-fa}
\nocite{*}
\bibliography{MyReferences}
}

\pagestyle{fancy}
\appendix
\chapter{جریان‌های کاری کامل فرآیند}
\newpage
\chapter{لیست کامل الگوریتم‌ها}
\newpage
\chapter{نتایج تجربی گسترده}
\newpage
\chapter{مصنوعات پیاده‌سازی و پکیج بازتولید}

\newpage

% چکیده و عنوان انگلیسی
\en-abstract{
(English summary of the thesis, required at the end by TMU guidelines)
}
\latinfirstPage

\end{document}
